% -*- mode : latex; coding: utf-8 -*-
\begin{KoreanAbstract}
  DRAM RowHammer 공격은 2014년 \cite{FlippingBits} 실험적으로 입증된 이후 컴퓨팅 시스템 보안에 
  중대한 위협으로 작용하였다. RowHammer 공격은 DRAM의 특정 wordline을 반복적으로 활성화 시켜 전자기파 
  간섭으로 인해 주변 다른 row의 축전기에서 전하가 빠져나가도록 유도하는 것으로, 반도체 공정이 점점 더 
  미세해짐에 따라 더 빈번하게 발생하고 있다 \cite{Revisiting}. RowHammer 공격을 막는 다양한 방식이 
  연구되었지만, DRAM 내부 구조가 공개되지 않았기 때문에 효율 측면에서 한계가 존재하였다. 
  본 논문에서는 DDR4 내부 구조가 공개되었을 때 제시된 RowHammer 방어전략들 중 가장 최신 방식인
   Grpahene \cite{Graphene} 과 Hydra \cite{Hydra} 에서 개선할 수 있는 부분들을 제시하였다. 
   RowHammer 공격이 최소 1K의 ACT 명령어에서 축전기의 전하를 충분히 빠져나가게 한다고 가정하고 
   Macsim \cite{MACSIM} 과 Ramulator \cite{RAMULATOR} (2 channel DDR4) 시뮬레이터를 사용해 
   2개의 RowHammer 공격 패턴과 SPEC CPU 2017에서 메모리 접근이 가장 많은 3개의 트레이스를 사용해 실험을 진행하였다.
    실험 결과 Graphene에서는 추가적인 REF 명령어의 수가 기존에 비해 공격 패턴에서는 50\%, SPEC 트레이스에서는 
    64.7\% 감소하였으며 IPC 손실을 65.7\% 감소시킬 수 있었다. 
    Hydra의 경우 추가적인 REF 명령어 수가 기존에 비해 공격 패턴에서는 91.4\%, SPEC 트레이스에서는 98.1\% 감소하였고 IPC 손실을 
    71.7\% 감소시킨 결과를 얻었다. 
\end{KoreanAbstract}
